\chapter{Fractures}
\index{Fractures}

\section*{Definition and Overview}
A fracture is a break in a bone, ranging from a fine crack to a complete break with the bone ends displaced. Fractures are among the most common injuries, most often caused by falls, direct blows, twisting, and road trauma, and they can occur in any bone of the body. The treatment goal is to restore the bone's alignment and hold it still while it heals, using a cast, splint, or surgery, and then to rehabilitate the limb so it returns to full function. Most fractures heal completely with good treatment. This chapter explains the types of fractures, their symptoms, and how they are diagnosed and treated.

\section*{Epidemiology}
Fractures affect around 1 in 15 people each year and are a leading reason for emergency department visits. The most common sites are the wrist (especially the distal radius), ankle, hip, and humerus. Fracture patterns differ with age: sports and road trauma dominate in the young, while falls on weakened, osteoporotic bone dominate in older adults, with hip fractures being the most serious. Wrist and hip fractures in older women are the hallmark of osteoporosis. Most fractures heal without lasting disability, but hip fractures in older adults are associated with significant mortality and loss of independence.

\section*{Aetiology and Pathophysiology}
\begin{itemize}
    \item \textbf{Mechanisms:} falls, direct blows, twisting injuries, road traffic accidents, and sport
    \item \textbf{Osteoporosis:} weakened bone fractures with minimal force, common in older adults, particularly women
    \item \textbf{Types of fracture:} simple (two pieces), comminuted (multiple pieces), displaced, open (through the skin), greenstick (incomplete, in children), and stress (overuse)
    \item \textbf{Healing process:} bleeding and inflammation at the break, formation of soft callus within weeks, hardening into new bone, and remodelling over months
    \item \textbf{Stability:} stable, well-aligned fractures heal in a cast; unstable or displaced fractures need internal fixation
    \item \textbf{Blood supply:} healing depends on blood flow; the scaphoid, femoral head, and tibia have relatively poor supplies and heal slowly
    \item \textbf{Growth plates:} in children, fractures involving the growth plate need careful follow-up because they can affect future growth
    \item \textbf{Factors affecting healing:} smoking, diabetes, poor nutrition, and infection delay healing
\end{itemize}

\section*{Clinical Features}
\begin{itemize}
    \item Pain at the site, worse with movement or weight bearing
    \item Swelling, bruising, and tenderness
    \item Deformity or abnormal angulation of the limb
    \item Inability to use the limb or bear weight
    \item A grating sensation or audible snap at the time of injury
    \item Open fractures: a wound over the fracture site
    \item In children: crying, refusal to move the limb, and tenderness
    \item Signs of nerve or vessel injury: numbness, tingling, or a pale, cold limb
\end{itemize}

\section*{Differential Diagnosis}
Injury symptoms need distinguishing from other injuries.
\begin{itemize}
    \item Sprains and ligament injuries, which are common in the ankle and knee
    \item Dislocations, in which the bone ends separate at a joint
    \item Contusions and muscle strains
    \item Tendon injuries, such as Achilles rupture
    \item Stress fractures, with gradual pain rather than an acute injury
    \item Pathological fractures through tumours or infection in weakened bone
    \item X-rays are the definitive diagnostic tool
\end{itemize}

\section*{When to Seek Medical Care}
Any injury with significant pain, swelling, deformity, or inability to use the limb should be assessed, with X-rays where indicated. Do not delay, because displaced fractures are easier to treat early.

\textbf{Contact your local emergency services for:}
\begin{itemize}
    \item Open fractures, with bone through the skin
    \item A pale, cold, numb, or blue limb
    \item Severe deformity, especially with a dislocated joint
    \item Fractures from high-energy trauma, falls from height, or road accidents
    \item A head injury, neck pain, or back pain with a fracture
    \item Inability to move the limb or loss of sensation
\end{itemize}

\section*{Diagnosis and Investigations}
\begin{itemize}
    \item \textbf{History and examination:} the mechanism of injury, and assessment of the limb, pulses, and sensation
    \item \textbf{X-rays:} two views of the injured bone, including the joint above and below
    \item \textbf{CT scan:} for complex fractures of the pelvis, heel, spine, and joints
    \item \textbf{MRI:} for stress fractures, ligament injury, and assessing bone viability
    \item \textbf{Follow-up imaging:} to confirm healing and alignment in the fracture clinic
\end{itemize}

\section*{Management}
\begin{itemize}
    \item \textbf{First aid:} immobilise the limb, apply ice, elevate, and seek care; do not try to straighten it
    \item \textbf{Reduction:} displaced fractures are realigned under anaesthesia or sedation
    \item \textbf{Non-operative treatment:} casts, splints, and braces for stable fractures
    \item \textbf{Surgical treatment:} plates, screws, rods, or external fixation for displaced, unstable, or open fractures
    \item \textbf{Pain management:} analgesics and elevation to reduce swelling
    \item \textbf{Weight-bearing restrictions:} as advised, varying with the bone and treatment
    \item \textbf{Physiotherapy:} essential to restore movement, strength, and function after healing
    \item \textbf{Bone health:} treat osteoporosis to prevent future fractures
    \item \textbf{Complication watch:} plaster care, wound checks, and review for pain or swelling
\end{itemize}

\section*{Complications}
\begin{itemize}
    \item Delayed union, non-union, and malunion
    \item Infection, especially in open fractures
    \item Nerve and blood vessel injury
    \item Compartment syndrome, a painful emergency from pressure in the limb
    \item Deep vein thrombosis after lower-limb fractures
    \item Stiffness, muscle wasting, and chronic pain
    \item Post-traumatic arthritis in fractures involving joints
\end{itemize}

\section*{Prognosis and Outcomes}
Most fractures heal completely within 6--12 weeks and return to full function with rehabilitation. Children's fractures heal fastest and remodel well. Surgical fixation allows earlier movement and better outcomes for displaced fractures. The main determinants of outcome are the fracture type, the quality of treatment, and the person's overall health: non-smokers with good nutrition and bone health heal reliably. With modern orthopaedic care, the majority of people recover fully and return to their previous activities.

\section*{Prevention and Self-Care}
\begin{itemize}
    \item Keep bones strong with calcium, vitamin D, and weight-bearing exercise
    \item Prevent falls with home safety, good lighting, and balance exercise
    \item Wear protective gear in sport and seatbelts in cars
    \item Build exercise loads gradually
    \item Do not smoke
    \item Follow cast, weight-bearing, and rehabilitation instructions
    \item Attend follow-up appointments
    \item Seek review for persistent pain, swelling, or change in the limb
\end{itemize}

\section*{Sources and Further Reading}
The following reputable sources provide further information for healthcare students and patients seeking reliable, current advice about fractures.
\begin{itemize}
    \item Healthdirect Australia. \textit{Fractures.} https://www.healthdirect.gov.au/fractures
    \item Australian Orthopaedic Association. \textit{Fracture care.} https://www.aoa.org.au/
    \item Healthy Bones Australia. \textit{Fracture prevention.} https://healthybonesaustralia.org.au/
    \item NHS. \textit{Fractures.} https://www.nhs.uk/conditions/fractures/
    \item Mayo Clinic. \textit{Broken bones.} https://www.mayoclinic.org/
    \item World Health Organization. \textit{Falls prevention.} https://www.who.int/
\end{itemize}
